% !TEX root = ../main.tex
\chapter{Pseudo-exam and second reading}\label{app:exam}
\begin{paracol}{2}

The long exercises of each chapter (more than about a quarter of an hour, more than three parts, or proofs that
need a new idea), moved here verbatim so that a first reading of the chapters is not held up by them. The
chapters keep short, targeted exercises, often the core of the long ones below. Use this appendix as a second
reading or, once the course is finished, as a cumulative self-test: retrieval spaced weeks after first study is
among the best-supported ways to make material stick (Cepeda et al.\ 2006; Adesope et al.\ 2017).
% once an exercise here gets a \begin{solution}, add \printsolutions at the end of this appendix

% ------------------------------------------------------------------
\section{Stochastic processes I (\cref{ch:sp1})}\label{exam:sp1}

\begin{exercise}[PS2 (2024), Problem 1: random walk with bias]\label{ch04:ex-ps2-1}
Let $X_n$ be i.i.d.\ steps with $\Prob(X_n=+1)=p$, $\Prob(X_n=0)=1-(p+q)$, $\Prob(X_n=-1)=q$, where $p,q>0$ and $p+q<1$
(symmetric if $p=q$, biased upwards if $p>q$), and $S_n=\sum_{i=1}^nX_i$, $S_0=0$.
\begin{enumerate}[(a)]
  \item Derive the moment generating function $M_{X_i}(t)$ of a step and the mean and variance of $X_i$.
  \item Derive the moment generating function $M_{S_n}(t)$ and the mean and variance of $S_n$.
  \item With $\mu$, $\sigma^2$ the mean and variance of a step, prove that $Z_n=(S_n-n\mu)/\sqrt{n\sigma^2}$ converges in
      distribution to $N(0,1)$.
  \item Model the intraday price of a stock as such a walk with 78 five-minute steps (the six and a half hours of NYSE
      trading), with $p=0.52$ and $q=0.47$. Which step size gives an expected daily price change of $+\$0.50$? With that
      step size, what is the standard deviation of the daily price change?
\end{enumerate}
Part (d) is \cref{ch04:ex-ps2-1-short} in the text.
\end{exercise}

\begin{exercise}[PS2 (2013), A-3: identify Markov processes]\label{ch04:ex-a3}
Which of the following are Markov processes?
\begin{enumerate}[(a)]
  \item The simple random walk $S_t$.
  \item $X_t=|S_t|$.
  \item $X_0=0$,
      $X_{t+1}=X_t+(Y_t-\tfrac1\lambda)$ with $Y_t$ i.i.d.\ exponential with parameter $\lambda$.
  \item $X_0=0$,
      $X_{t+1}=X_t+Z_tZ_{t-1}\cdots Z_0$ with $Z_t$ i.i.d.\ log-normal.
  \item $X_0=0$, $X_{t+1}=X_t+W_t+W_{t-1}+\dots+W_0$ with
      $W_t$ i.i.d.\ normal.
\end{enumerate}
Parts (b) and (e) are \cref{ch04:ex-a3-short} in the text.
\end{exercise}

\begin{exercise}[Stopping times (not from the course)]\label{ch04:ex-stoptimes}
Let $\sigma,\tau$ be stopping times. Show that $\sigma\wedge\tau$, $\sigma\vee\tau$ and $\tau+1$ are stopping times, but that
$\tau-1$ in general is not. Deduce that ``one step after the first peak'' is a stopping time and ``the first peak'' is
not (\cref{ch04:ex-stop}). The part on $\tau\pm1$ is \cref{ch04:ex-stoptimes-short} in the text.
\end{exercise}

\begin{exercise}[Biased ruin by first-step analysis (not from the course)]\label{ch04:ex-biasedfsa}
For the biased walk let $f(k)=\Prob(\text{hit }A\text{ before }-B\mid S_0=k)$. Show that
$f(k)=pf(k+1)+qf(k-1)$, solve this linear recursion (characteristic roots $1$ and $q/p$) with $f(A)=1$, $f(-B)=0$, and
recover \cref{ch04:prop-biased}. Then use Wald's identity to show that the expected duration tends to $AB$ as
$p\to\frac12$. A checking version is \cref{ch04:ex-biasedfsa-short} in the text.
\end{exercise}

% ------------------------------------------------------------------
\section{Regression (\cref{ch:regression})}\label{exam:regression}

\begin{exercise}[not from the course: optimal design of a line fit]\label{ch05:exam-1}
For $y=\beta_0+\beta_1x+\varepsilon$ under Gauss--Markov with $n$ even and $x_i\in[-1,1]$ at your choice:
\begin{enumerate}[(a)]
  \item derive
      $\hat\beta_1=S_{xy}/S_{xx}$ and $\Var(\hat\beta_1)=\sigma^2/S_{xx}$ with $S_{xx}=\sum_i(x_i-\bar x)^2$;
  \item find the design minimising
      $\Var(\hat\beta_1)$ and compute the leverages $H_{ii}$ for it;
  \item explain why this design cannot detect curvature, and how the
      curvature test of \texttt{residualPlots()} would behave.
\end{enumerate}
Its core is \cref{ch05:exam-1-short} in the text.
\end{exercise}

\begin{exercise}[not from the course: the Chow test]\label{ch05:exam-2}
\leavevmode
\begin{enumerate}[(a)]
  \item Show that the residual sum of squares of the block-diagonal design $\bigl(\begin{smallmatrix}X_A&0\\0&X_B\end{smallmatrix}\bigr)$ is
      $\text{RSS}_A+\text{RSS}_B$.
  \item Show that $\bigl(\begin{smallmatrix}X_A&0\\X_B&X_B\end{smallmatrix}\bigr)$ has the same column space, and relate the two
      parameter vectors.
  \item For GE, compute $\hat\alpha_B$, $\hat\beta_B$ and the $F$ statistic from the numbers in the case study of
      \cref{ch05:sec-chow}. Why can you not obtain the standard error of $\hat\beta_B$ from that table alone?
\end{enumerate}
Its core is \cref{ch05:exam-2-short} in the text.
\end{exercise}

\begin{exercise}[PS3 (2013), Problem 1(d)--(e): invariance and leverage]\label{ch05:ex-ps3-1de}
(d) Prove that $H$ is unchanged if $X$ is replaced by $X'=XG$ for any invertible $p\times p$ matrix $G$.
(e) Let $X$ be $n\times(p+1)$ with a first column of ones and $p$ covariates, and let $G$ be the identity except for the first
row, which is $(1,-\bar x_1,\dots,-\bar x_p)$, with $\bar x_j=\frac1n\sum_ix_{i,j}$. By (d) the models with $X$ and $X'=XG$ have the
same fitted values and residuals.
\begin{itemize}
  \item If $\bm\beta=(\beta_0,\dots,\beta_p)\T$ is the parameter for $X$, show that $\bm\beta'=G^{-1}\bm\beta$ is the parameter for $X'$; find
        $G^{-1}$ and the elements of $\bm\beta'$ explicitly.
  \item Show that $X'^{\mathsf T}X'=\bigl(\begin{smallmatrix}n&\bm 0_p\T\\ \bm 0_p&\mathcal X\T\mathcal X\end{smallmatrix}\bigr)$, where $\mathcal X$ has entries $x_{i,j}-\bar x_j$.
  \item Prove $H_{ij}=\frac1n+(\bm x_i-\bar{\bm x})\T[\mathcal X\T\mathcal X]^{-1}(\bm x_j-\bar{\bm x})$, where $\bm x_i=(x_{i,1},\dots,x_{i,p})\T$. Hence the
        leverage $H_{ii}$ increases with the (squared) Mahalanobis distance between $\bm x_i$ and $\bar{\bm x}$.
\end{itemize}
Its core is \cref{ch05:ex-ps3-1de-short} in the text.
\end{exercise}

\begin{exercise}[not from the course: ridge as augmented least squares]\label{ch05:exam-3}
\leavevmode
\begin{enumerate}[(a)]
  \item Show that $\hat{\bm\beta}^{\text{ridge}}$ is the OLS estimate for the augmented data $\tilde{\bm y}=(\bm y,\bm 0_p)$,
      $\tilde X=\bigl(\begin{smallmatrix}X\\ \sqrt\lambda I_p\end{smallmatrix}\bigr)$.
  \item Using the SVD, show that $\|\hat{\bm\beta}^{\text{ridge}}\|$ decreases in $\lambda$ and compute
      $\Cov(\hat{\bm\beta}^{\text{ridge}})$ under Gauss--Markov.
  \item Show that ridge has smaller variance than OLS in every direction:
      $\Cov(\hat{\bm\beta})-\Cov(\hat{\bm\beta}^{\text{ridge}})\succeq0$.
\end{enumerate}
Its core is \cref{ch05:exam-3-short} in the text.
\end{exercise}

\begin{exercise}[PS3 (2013), Problem 2(b)--(d): case-deletion influence]\label{ch05:ex-ps3-2bd}
(b) Using Sherman--Morrison--Woodbury, show that the least-squares estimate with case $i$ deleted is
$\hat{\bm\beta}_{(i)}=\hat{\bm\beta}-(X\T X)^{-1}\bm x_i\hat\varepsilon_i/(1-H_{ii})$, where $\bm x_i\T$ is the $i$-th row of $X$ and $\hat\varepsilon_i=y_i-\bm x_i\T\hat{\bm\beta}$.
(c) Cook's distance is $CD_i=\|\hat{\bm y}-\hat{\bm y}_{(i)}\|^2/(p\hat\sigma^2)$ with $\hat{\bm y}_{(i)}=X\hat{\bm\beta}_{(i)}$. Show that
$CD_i=\dfrac{\hat\varepsilon_i^2}{p\hat\sigma^2}\cdot\dfrac{H_{ii}}{(1-H_{ii})^2}$.
(d) Let $\hat\sigma^2_{(i)}$ be the unbiased estimate of $\sigma^2$ with case $i$ deleted. Show that
$\hat\sigma^2_{(i)}=\hat\sigma^2+\frac{1}{n-p-1}\bigl(\hat\sigma^2-\frac{\hat\varepsilon_i^2}{1-H_{ii}}\bigr)$.
\end{exercise}

\begin{exercise}[PS3 (2013), Problem 3(c)--(d): sequential ANOVA]\label{ch05:ex-ps3-3cd}
Let $A=\bigl(\begin{smallmatrix}Q\T\\ W\T\end{smallmatrix}\bigr)$ be orthogonal as in the proof of \cref{ch05:thm-normal} and $\bm z=A\bm y=(z_1,\dots,z_n)\T$.
(c) Prove that $\bm y\T\bm y=\sum_{i=1}^nz_i^2$, $\hat{\bm y}\T\hat{\bm y}=\sum_{i=1}^pz_i^2$, $\hat{\bm\varepsilon}\T\hat{\bm\varepsilon}=\sum_{i=p+1}^nz_i^2$, and for the constrained model
$\hat{\bm y}_I\T\hat{\bm y}_I=\sum_{i=1}^kz_i^2$, $\hat{\bm\varepsilon}_I\T\hat{\bm\varepsilon}_I=\sum_{i=k+1}^nz_i^2$.
(d) Using $\bm z_Q\sim N_p(R\bm\beta,\sigma^2I_p)$ independent of $\bm z_W\sim N_{n-p}(\bm 0,\sigma^2I_{n-p})$, deduce that
$SS_{\text{ERROR}}=\hat{\bm\varepsilon}\T\hat{\bm\varepsilon}\sim\sigma^2\chi^2_{n-p}$; that under $H_0$
$SS_{\text{REG}(k+1,\dots,p\mid1,\dots,k)}=\hat{\bm y}\T\hat{\bm y}-\hat{\bm y}_I\T\hat{\bm y}_I=\sum_{i=k+1}^pz_i^2\sim\sigma^2\chi^2_{p-k}$, independent of
$SS_{\text{ERROR}}$; and that $\hat F=\dfrac{SS_{\text{REG}}/(p-k)}{SS_{\text{ERROR}}/(n-p)}\sim F_{p-k,\,n-p}$. Arrange the results in the ANOVA
table with rows ``regression on $1,\dots,k$'' ($\hat{\bm y}_I\T\hat{\bm y}_I$, $k$ d.f.), ``regression on $k+1,\dots,p$ adjusting for
$1,\dots,k$'' ($p-k$ d.f., mean square $MS_0$, $\hat F=MS_0/MS_{\text{Error}}$), ``error'' ($n-p$ d.f.) and ``total'' ($\bm y\T\bm y$,
$n$ d.f.).
\end{exercise}

% ------------------------------------------------------------------
\section{Linear rates products (\cref{ch:rates})}\label{exam:ch06}

\begin{exercise}[not from the course]\label{ch06:exam-1}
Using the quotes of \cref{ch06:ex-boot}:

(a) bootstrap $Z_1,Z_2,Z_3$ by hand;

(b) show that CDF gives $Z_4=\sqrt{Z_3Z_5}$, write the scalar
equation for $Z_5$ and solve it numerically;

(c) compute the one-year forward rates $f(k-1,k)$,
$k=1,\dots,5$, and the 4-year par rate;

(d) prove \cref{ch06:prop-cdf} directly from the daily compounding
definition, and show that the forward curve jumps at a knot unless
the two adjacent segments have equal slopes of $\ln Z$.
Its core is \cref{ch06:exam-1-short} in the text.
\end{exercise}

\begin{exercise}[not from the course]\label{ch06:exam-2}
\leavevmode
\begin{enumerate}[(a)]
  \item Show that the gamma of slide 31 equals
      $\epsilon\,\partial^2\mathrm{PV}/\partial r^2+O(\epsilon^3)$.
  \item For a 10-year zero-coupon bond with continuously compounded
      rate 4\% and notional \$100M, compute the PV01 and this gamma for
      $\epsilon=1$\,bp.
  \item How many revaluations does a full central-difference
      cross-gamma matrix need for $N$ curve inputs?
  \item Why does a Monte Carlo pricer need a larger $\epsilon$
      (\cref{ch:markets})?
\end{enumerate}
Its core is \cref{ch06:exam-2-short} in the text.
\end{exercise}

% ------------------------------------------------------------------
\section{PCA and factor models (\cref{ch:pca})}\label{exam:ch07}

\begin{exercise}[PS2 (2024), Problem 3: principal components, special case]\label{ch07:ex-ps2}
A portfolio holds $n$ assets with fixed weights $\bm w=(w_1,\dots,w_n)\T$; over a given horizon the assets have
rates of return $\bm X=(X_1,\dots,X_n)\T$ with mean $\bm\mu$ and covariance $\Sigma=(\sigma_{ij})$, and the portfolio return
is $Y=\bm w\T\bm X$.
\begin{enumerate}[(a)]
  \item Show that $\E[Y]=\bm w\T\bm\mu$ and $\Var[Y]=\bm w\T\Sigma\bm w$.
  \item Prove that $\Sigma$ is positive semi-definite, for any random variables $X_1,\dots,X_n$.
  \item Suppose all entries of $\Sigma$ are positive. Prove that the largest eigenvalue $\lambda_{\max}$ has multiplicity one and
      an eigenvector $\bm v$ that can be chosen with all components positive.
  \item Write the first principal component variable $PC1$ in terms of $\bm X$, $\bm\mu$ and $\bm v$. What are its mean and
      variance?
\end{enumerate}
Its core is \cref{ch07:ex-ps2-short} in the text.
\end{exercise}

\begin{exercise}[PS3 (2024), Problem 4; also PS7 (2013), Problem 3: PCA of yield changes]\label{ch07:ex-ps3}
Daily changes of the 9 FRED constant-maturity Treasury yields over 2001--2005 are analysed twice, with the sample
covariance matrix and with the sample correlation matrix. The standard deviations of the PC variables are
\begin{center}\small
\begin{tabular}{@{}l rrrrrrrrr@{}}
\toprule
 & 1 & 2 & 3 & 4 & 5 & 6 & 7 & 8 & 9\\
\midrule
covariance & 0.1618 & 0.0532 & 0.0299 & 0.0192 & 0.0135 & 0.0112 & 0.0096 & 0.0090 & 0.0081\\
correlation & 2.6338 & 1.1892 & 0.5652 & 0.3971 & 0.2694 & 0.2080 & 0.1485 & 0.1366 & 0.1225\\
\bottomrule
\end{tabular}
\end{center}
(proportions of variance 84.9, 9.2, 2.9\,\% and 77.1, 15.7, 3.5\,\% for the first three) and the first three loading
vectors are
\begin{center}\small
\begin{tabular}{@{}l rrr c rrr@{}}
\toprule
 & \multicolumn{3}{c}{covariance PCA} & & \multicolumn{3}{c}{correlation PCA}\\
\cmidrule(lr){2-4}\cmidrule(l){6-8}
Tenor & PC1 & PC2 & PC3 & & PC1 & PC2 & PC3\\
\midrule
3M  & 0.1088 & 0.5225 & 0.5434 & & $-0.2083$ & 0.6293 & 0.5831\\
6M  & 0.1581 & 0.4817 & 0.2508 & & $-0.2830$ & 0.5154 & $-0.0654$\\
1Y  & 0.2531 & 0.4107 & $-0.0514$ & & $-0.3369$ & 0.2698 & $-0.4066$\\
2Y  & 0.3905 & 0.2071 & $-0.4750$ & & $-0.3625$ & 0.0064 & $-0.4005$\\
3Y  & 0.4197 & 0.0873 & $-0.4022$ & & $-0.3664$ & $-0.0665$ & $-0.2719$\\
5Y  & 0.4206 & $-0.1082$ & $-0.0311$ & & $-0.3670$ & $-0.1619$ & 0.0036\\
7Y  & 0.4009 & $-0.2289$ & 0.1475 & & $-0.3608$ & $-0.2298$ & 0.1435\\
10Y & 0.3697 & $-0.2837$ & 0.2710 & & $-0.3525$ & $-0.2656$ & 0.2554\\
20Y & 0.3101 & $-0.3621$ & 0.3944 & & $-0.3290$ & $-0.3338$ & 0.4125\\
\bottomrule
\end{tabular}
\end{center}
(a) PC1 loadings are all positive in the covariance case and all negative in the correlation case. Is the
difference in sign meaningful? (Does the eigen-decomposition change if an eigenvector is multiplied by $-1$?)
(b) The PC1 loadings range more widely in the covariance case; the loading on the least variable change (3M) is
larger in magnitude in the correlation case, and those on the most variable changes are smaller. Explain why,
recalling that PC1 is the unit-norm weighted combination with the largest variance and that in the correlation
case each yield change is scaled to unit variance.
(c) Interpret the first three PCs of the correlation case and compare with the covariance case.
(d) If the goal is to model the dynamics of the whole term structure across all tenors, argue why correlation PCA
might be preferred.
(e) (2024 only) Compute the total variance of the correlation-case PCs. Is it an integer? Explain.
\end{exercise}

\begin{exercise}[PS7 (2013), Problem 1: PCA of a bivariate vector]\label{ch07:ex-ps7-1}
Let $\bm X=(X_1,X_2)\T$ have mean $\bm\alpha=(\alpha_1,\alpha_2)\T$ and covariance
$\Sigma=\bigl(\begin{smallmatrix}\sigma_1^2&\rho\sigma_1\sigma_2\\ \rho\sigma_1\sigma_2&\sigma_2^2\end{smallmatrix}\bigr)$.
\begin{enumerate}[(a)]
  \item Compute the eigenvalues $\lambda_1\ge\lambda_2\ge0$ of $\Sigma$.
  \item Compute orthonormal eigenvectors $\bm\gamma_1,\bm\gamma_2$.
  \item Verify $\Sigma=\sum_i\lambda_i\bm\gamma_i\bm\gamma_i\T$.
  \item For $p_i=\bm\gamma_i\T(\bm X-\bm\alpha)$ prove $\E[p_i]=0$, $\Var(p_i)=\lambda_i$ and $\Cov(p_1,p_2)=0$.
  \item Are the eigenvectors in (b) unique? If not, how are different solutions related?
\end{enumerate}
Its core is \cref{ch07:ex-ps7-1-short} in the text.
\end{exercise}

\begin{exercise}[PS7 (2013), Problem 4: dynamics of the PC scores]\label{ch07:ex-ps7-4}
A VAR(3) with constant (\cref{ch:timeseries}) was fitted to the scores of the first three PCs of the
correlation-matrix analysis of \cref{ch07:ex-ps3} ($T=1245$ days). Estimates and $t$-values:
\begin{center}\small
\begin{tabular}{@{}l rr c rr c rr@{}}
\toprule
 & \multicolumn{2}{c}{eq.\ Comp.1} & & \multicolumn{2}{c}{eq.\ Comp.2} & & \multicolumn{2}{c}{eq.\ Comp.3}\\
\cmidrule(lr){2-3}\cmidrule(lr){5-6}\cmidrule(l){8-9}
Regressor & est. & $t$ & & est. & $t$ & & est. & $t$\\
\midrule
Comp.1, lag 1 & 0.0527 & 1.86 & & $-0.0203$ & $-1.65$ & & $-0.0029$ & $-0.48$\\
Comp.2, lag 1 & $-0.1340$ & $-2.08$ & & 0.1382 & 4.93 & & $-0.0143$ & $-1.03$\\
Comp.3, lag 1 & $-0.1322$ & $-1.00$ & & $-0.0833$ & $-1.45$ & & 0.0163 & 0.58\\
Comp.1, lag 2 & $-0.0530$ & $-1.87$ & & $-0.0072$ & $-0.59$ & & $-0.0135$ & $-2.23$\\
Comp.2, lag 2 & $-0.0542$ & $-0.84$ & & 0.0370 & 1.32 & & $-0.0037$ & $-0.27$\\
Comp.3, lag 2 & 0.1825 & 1.39 & & $-0.0195$ & $-0.34$ & & $-0.0523$ & $-1.86$\\
Comp.1, lag 3 & $-0.0169$ & $-0.60$ & & $-0.0085$ & $-0.69$ & & $-0.0019$ & $-0.32$\\
Comp.2, lag 3 & 0.0130 & 0.21 & & $-0.0403$ & $-1.46$ & & $-0.0030$ & $-0.22$\\
Comp.3, lag 3 & 0.0830 & 0.63 & & 0.0148 & 0.26 & & $-0.0712$ & $-2.53$\\
constant & $-0.0114$ & $-0.15$ & & 0.0121 & 0.38 & & 0.0002 & 0.01\\
\midrule
$R^2$; $F$-test $p$-value & 0.013 & 0.054 & & 0.029 & $4\cdot10^{-5}$ & & 0.013 & 0.059\\
\bottomrule
\end{tabular}
\end{center}
Residual standard errors are 2.60, 1.13, 0.56; residual correlations are below $0.025$ in absolute value; the
largest root of the characteristic polynomial has modulus 0.438.
(a)--(c) Interpret the estimated equations for Comp.1, Comp.2 and Comp.3.
(d) For each equation, is there evidence of mean reversion or of momentum in the PC variable? Explain.
\end{exercise}

\begin{exercise}[not from the course]\label{ch07:ex-fly}
Using the covariance-PCA loadings of \cref{ch07:ex-ps3}:
\begin{enumerate}[(a)]
  \item find the PC1-neutral 2s10s steepener and check
      that its exposure to PC2 is non-zero;
  \item find the PC1/PC2-neutral 2s5s10s fly (buy \$1 dv01 of 5y) and compute its
      exposure to PC3;
  \item prove the variance formula for the PC1-hedged portfolio $\bm V'=\bm V-(\bm V\T\bm\gamma_1)\bm\gamma_1$ and show
      that it never exceeds $\Var(\bm V\T\bm r)$.
\end{enumerate}
\end{exercise}
